Esta gráfica resume toda la comparación en una imagen. Cada fila es un rival. El círculo hueco es el paper; el relleno, la réplica con Docker; la rayita, la réplica sin Docker. La raya vertical es Autores: a la izquierda, el rival pierde; a la derecha, gana. La escala es logarítmica.

Si hueco y relleno están juntos, la réplica clavó al paper. Pasa con Libstdc++, con Williams y con Anderson en casi todos los escenarios.

La excepción salta a la vista: Folly, en amarillo. En el paper está muy a la izquierda, rindiendo un 5 por ciento de Autores. Aquí está cerca de la raya, o incluso a la derecha. Es la mayor discrepancia de la réplica.

Y fijaos en la rayita sin Docker: cae casi siempre sobre el círculo relleno. Docker no cambia nada, lo veremos en la parte 4.
